\chapter{Flutter: reportar, moderar y administrar}
\label{ch:flutter_admin}

\section{Objetivo}

Implementar en el cliente Flutter toda la interfaz administrativa y de moderacion:

\begin{itemize}[leftmargin=*]
  \item Boton de reportar en publicaciones y canales.
  \item Menu contextual en cada publicacion con Reportar y Eliminar.
  \item Pantalla de administracion con tres pestanas: Reportes, Usuarios y Auditoria.
  \item Detalle del reporte con votacion.
  \item Gestion de roles y baneo de usuarios.
  \item Consulta del log de auditoria.
\end{itemize}

\section{Contexto}

El backend de la Fase~\ref{ch:backend_reportes} ya exponia todos los endpoints. El cliente no tenia ninguna interfaz que los consumiera. Se dividio el trabajo en cinco sub-fases: D2 (boton de reportar), D3 (lista de reportes), D4 (panel completo), D5 (moderacion directa).

\section{Desarrollo}

\subsection{Modelos nuevos}

\begin{itemize}[leftmargin=*]
  \item \texttt{Report}: mapea la respuesta de \texttt{/api/admin/reports}.
  \item \texttt{AdminUser}: mapea \texttt{/api/admin/users}.
  \item \texttt{AuditLog}: mapea \texttt{/api/admin/audit}.
\end{itemize}

\subsection{Providers nuevos}

\begin{itemize}[leftmargin=*]
  \item \texttt{ReportProvider}: crear reportes desde cualquier usuario.
  \item \texttt{AdminProvider}: carga de reportes, usuarios y auditoria, ademas de votar, cambiar rol y banear.
  \item \texttt{PostProvider.delete()}: borra publicaciones y actualiza el estado local.
\end{itemize}

\subsection{Wigets reutilizables}

\subsubsection{\texttt{ReportDialog}}

Modal reutilizable con:

\begin{itemize}[leftmargin=*]
  \item Radio buttons con las cinco razones (\texttt{SPAM}, \texttt{INAPPROPRIATE}, \texttt{MISINFORMATION}, \texttt{HARASSMENT}, \texttt{OTHER}).
  \item Campo de detalles opcional.
  \item Validacion y feedback con SnackBar.
\end{itemize}

Se expone como metodo estatico:

\begin{lstlisting}[language=Java]
static Future<void> show(BuildContext context, {
  required String targetType,   // 'POST' | 'CHANNEL'
  required String targetId,
  required String targetLabel,
});
\end{lstlisting}

\subsection{Menu contextual en publicaciones}

Se reemplazo el icono de reportar por un \texttt{PopupMenuButton} que muestra las opciones segun los permisos del usuario actual.

\begin{lstlisting}[language=Java]
final isOwn = currentUserId != null && post.authorId == currentUserId;
final canDelete = isOwn || canModerateAny;
final canReport = !isOwn;

PopupMenuButton<String>(
  onSelected: (value) {
    if (value == 'report') ReportDialog.show(...);
    else if (value == 'delete') onDelete();
  },
  itemBuilder: (_) => [
    if (canReport) PopupMenuItem(value: 'report', child: ...),
    if (canDelete) PopupMenuItem(value: 'delete', child: ...),
  ],
);
\end{lstlisting}

\subsection{Pantalla \texttt{AdminScreen}}

Contiene un \texttt{TabBar} con tres sub-pestanas:

\begin{itemize}[leftmargin=*]
  \item \texttt{AdminReportsTab}: lista con filtros por estado.
  \item \texttt{AdminUsersTab}: listado con buscador y acciones.
  \item \texttt{AdminAuditTab}: bitacora con iconos por tipo de accion.
\end{itemize}

Solo es visible para ADMIN y SUPERADMIN (verificado con \texttt{auth.isAdminOrSuper}).

\subsection{\texttt{AdminReportsTab}}

Filtros por estado:

\begin{lstlisting}[language=Java]
static const _filters = [
  ('PENDING', 'Pendientes'),
  ('RESOLVED_KEEP', 'Conservados'),
  ('RESOLVED_DELETE', 'Eliminados'),
  ('ALL', 'Todos'),
];
\end{lstlisting}

Cada tarjeta muestra:

\begin{itemize}[leftmargin=*]
  \item Chip de tipo (Post/Canal).
  \item Chip de razon.
  \item Chip de estado con color.
  \item Preview del contenido reportado.
  \item Contador de votos \texttt{0/3}, \texttt{1/3}, etc.
  \item Indicador de voto propio si el usuario ya voto.
\end{itemize}

\subsection{\texttt{ReportDetailScreen}}

Pantalla de detalle con:

\begin{itemize}[leftmargin=*]
  \item Contenido reportado y su ubicacion.
  \item Motivo y detalles del reporte.
  \item Informacion del reporter.
  \item Contador de votos con barra de progreso.
  \item Botones de voto (Conservar / Eliminar).
  \item Estado de resuelto si aplica.
\end{itemize}

Antes de votar, se pide un comentario opcional en un dialogo.

\subsection{\texttt{AdminUsersTab}}

Lista con:

\begin{itemize}[leftmargin=*]
  \item Avatar con inicial.
  \item Nombre, email y badge de rol con color.
  \item Badge ``Tu'' en el usuario actual.
  \item Badge ``BANEADO'' si esta deshabilitado.
  \item Boton de tres puntos que abre un bottom sheet con las acciones.
\end{itemize}

\subsection{Bottom sheet de acciones}

Para ADMIN muestra solo promociones (Estudiante, Profesor, Administrador).
Para SUPERADMIN agrega:

\begin{itemize}[leftmargin=*]
  \item Opcion de nombrar SUPERADMIN.
  \item Boton de Banear o Reactivar.
\end{itemize}

Cada accion pide confirmacion en un dialogo antes de ejecutarse.

\subsection{\texttt{AdminAuditTab}}

Muestra el log con:

\begin{itemize}[leftmargin=*]
  \item Icono y color segun la accion (\texttt{VOTE\_KEEP}, \texttt{VOTE\_DELETE}, \texttt{CHANGE\_ROLE}, \texttt{BAN\_USER}, \texttt{UNBAN\_USER}, \texttt{DELETE\_POST}).
  \item Nombre del admin.
  \item Detalles.
  \item Fecha y hora.
\end{itemize}

Un banner superior indica si esta viendo todo el log (SUPERADMIN) o solo el propio (ADMIN).

\subsection{Cambios a archivos previos}

\begin{itemize}[leftmargin=*]
  \item \texttt{home\_shell.dart}: se agrego la pestana Admin condicional segun el rol.
  \item \texttt{channel\_detail\_screen.dart}: se agrego el menu contextual por publicacion.
  \item \texttt{post\_provider.dart}: se agrego el metodo \texttt{delete}.
  \item \texttt{main.dart}: se registraron \texttt{ReportProvider} y \texttt{AdminProvider}.
\end{itemize}

\section{Verificacion}

\begin{itemize}[leftmargin=*]
  \item Un estudiante puede reportar publicaciones y canales.
  \item Un administrador ve la pestana Admin.
  \item El SUPERADMIN ve tres sub-pestanas y puede votar, gestionar usuarios y consultar el log completo.
  \item Un ADMIN solo ve sus propias acciones en el log.
  \item El dueno del canal puede borrar publicaciones ajenas sin votacion.
\end{itemize}

\section{Resultado}

\begin{exito}[Estado al cierre]
\begin{itemize}[leftmargin=*]
  \item Interfaz completa de reportes y moderacion.
  \item Gestion de usuarios con jerarquia de roles.
  \item Auditoria visible segun el rol.
  \item Moderacion directa por dueno del canal.
\end{itemize}
\end{exito}